\chapter{Ambiguity and robust choice}
\label{sec:p17-ambiguity-and-robust-choice}

Ellsberg's examples ask whether one additive probability assessment
and one consequence-utility scale can represent a person's choices
across specified bets. The contradiction is a precise statement about
that representation. A separate question is whether the choices are
predictable, defensible, or common in a population.
The original article develops those questions through hypothetical
comparisons and reported conversations
\cite[pp.~643--669]{Ellsberg1961Ambiguity}.

After the original comparisons, this chapter gives a finite
set-of-priors maxmin method. Its source is the actual
Gilboa--Schmeidler lottery-act formulation
\cite[pp.~141--145]{GilboaSchmeidler1989Maxmin}.
Here the utility and prior set are supplied inputs. The finite
calculations and forward properties below do not constitute the
general axiomatic representation theorem of that paper.

\section{Acts, a single prior, and the sure-thing condition}

\begin{hbdefinition}{def:p17-ambiguity-acts}{Consequences, lottery acts, and one-prior evaluation}
Let \(S\) and \(Z\) be nonempty finite sets of states and consequences.
An act \(f:S\to\Delta(Z)\) specifies a lottery over consequences at
each state. A deterministic act is identified with its degenerate
lotteries. Fix a utility \(u:Z\to\mathbb R\), and define
\[
 x_f(s)=\sum_{z\in Z}f(s)(z)u(z),\qquad
 E_p(f)=\sum_{s\in S}p(s)x_f(s),\quad p\in\Delta(S).
\]
Thus a prior is a nonnegative vector of total mass one and is
additive on events. Subjective expected utility with one prior
orders all acts by this same \(E_p\), with utility independent
of the state and of which bet is being compared.
The utility need not be linear in money.
The finite lottery domain and its mixture probabilities are
those of \cref{def:p17-finite-lottery,thm:p17-lottery-reduction}.
This consequence-utility model is distinct from the additive
monetary-bet assumptions of \cref{def:p17-bet-axioms}.

A constant act gives the same lottery at every state.
For a weak preference, \(f\succ g\) means
\(f\succeq g\) and not \(g\succeq f\); \(f\sim g\)
means both weak comparisons hold.
For \(E\subseteq S\), the pasted act \(f_Eh\) equals \(f\)
on \(E\) and \(h\) on \(S\setminus E\).
A complete transitive preference satisfies the
\emph{sure-thing condition} when, for all these acts and events,
\[
       f_Eh\succeq g_Eh
                  \quad\Longleftrightarrow\quad
       f_Eh'\succeq g_Eh'.
\]
Only consequences on which the two compared acts agree have
changed. This is the finite pasting form of the P2 comparison
used by Ellsberg \cite[pp.~648--649]{Ellsberg1961Ambiguity}.
It is distinct from independence of probabilistic events and
from the mixture-independence axiom defined below.
\index{ambiguity}\index{sure-thing principle}
\end{hbdefinition}

\begin{hblemma}{thm:p17-ambiguity-seu-benchmark}{proved=annex:p17.ambiguity.seu.benchmark}{The additive benchmark and its exact restrictions}
Every one-prior expected-utility preference satisfies the
sure-thing condition. If a bet \(b_E\) pays consequence \(H\)
on \(E\) and \(L\) otherwise, with \(u(H)>u(L)\), then
\[
              E_p(b_E)=u(L)+(u(H)-u(L))p(E).
\]
Therefore the ranking of two such bets is precisely the
ranking of their event probabilities. Complementation reverses
that ranking; adjoining the same disjoint event preserves it.
These conclusions use the common prize pair and common utility
scale. An indifference between two bets by itself does not
establish that all the person's other choices have this representation.
\end{hblemma}

\section{Ellsberg's original urn comparisons}

\begin{hbdefinition}{def:p17-ellsberg-two-urns}{Two urns and four offered comparisons}
Urn I contains 100 red and black balls in an unknown proportion,
possibly from zero to 100 red balls. Urn II is known to contain
50 red and 50 black balls. A bet pays \(H=\$100\) if the
selected color is drawn uniformly from the selected urn and
\(L=\$0\) otherwise. Write \(R_{\rm I},B_{\rm I},
R_{\rm II},B_{\rm II}\) for the corresponding four bets.
Ellsberg offers the comparisons
\[
 R_{\rm I}\ \hbox{versus}\ B_{\rm I},\quad
 R_{\rm II}\ \hbox{versus}\ B_{\rm II},\quad
 R_{\rm I}\ \hbox{versus}\ R_{\rm II},\quad
 B_{\rm I}\ \hbox{versus}\ B_{\rm II}.
\]
His principal pattern combines within-urn indifference with
preference for the known urn for either color. He also records
the reversed cross-urn pattern and complete indifference
\cite[pp.~650--653]{Ellsberg1961Ambiguity}.
These are separate offered comparisons; freely choosing an urn
and then randomizing the color is an additional available act.
\end{hbdefinition}

\begin{hbtheorem}{thm:p17-ellsberg-two-urns}{proved=annex:p17.ellsberg.two.urns}{The complementary-bet contradiction}
For the common prizes with \(u(H)>u(L)\), no one additive
subjective prior represents
\[
 R_{\rm II}\succ R_{\rm I}
                    \quad\hbox{and}\quad
 B_{\rm II}\succ B_{\rm I}.
\]
The reversed pair of strict preferences is equally incompatible.
More generally, weak preference for one urn for both colors
with at least one strict preference is impossible in this
representation. This conclusion does not require the observer
to infer a numerical probability from the within-urn indifferences.
If all four bets are indifferent, probabilities one half for
each color in each urn do represent those comparisons.
\end{hbtheorem}

\begin{hbdefinition}{def:p17-ellsberg-marked-urn}{The marked-ball and Arrow variants}
In Ellsberg's marked variant, combine the urns into one urn:
100 balls bear label I with unknown red/black split, and the
50 red and 50 black known-urn balls bear label II.
Draw one of the 200 balls uniformly.
In state order
\((R_{\rm I},B_{\rm I},R_{\rm II},B_{\rm II})\), the six
acts below list their consequence in each state:
\[
\begin{array}{c|cccc}
 M_1&H&L&L&L\\
 M_2&L&H&L&L\\
 M_3&L&L&H&L\\
 M_4&L&L&L&H\\
 M_5&H&H&L&L\\
 M_6&L&L&H&H
\end{array}
\]
The assumptions are \(M_1\sim M_2\), \(M_3\sim M_4\),
and \(M_5\sim M_6\).
Arrow's separate four-act variant on the same state space is
\[
\begin{array}{c|cccc}
 A_1&H&H&L&L\\
 A_2&H&L&H&L\\
 A_3&L&H&L&H\\
 A_4&L&L&H&H
\end{array}
\]
with \(A_1\sim A_4\), \(A_2\sim A_3\), and
\(A_1\succ A_2\).
These are the original additional comparisons
\cite[p.~652 and p.~654, footnote~4]{Ellsberg1961Ambiguity};
the marked draw is a different experiment from choosing one
of the original two urns.
\end{hbdefinition}

\begin{hbtheorem}{thm:p17-ellsberg-marked-arrow}{proved=annex:p17.ellsberg.marked.arrow}{The full pasting arguments in the two variants}
A complete transitive preference satisfying the sure-thing
condition and the three marked-urn indifferences must have
\[
                         M_1\sim M_2\sim M_3\sim M_4.
\]
In Arrow's variant, the stated assumptions force
\(A_4\succ A_3\) by transitivity, while the sure-thing
condition forces \(A_3\succ A_4\). Thus they cannot all hold.
These proofs require no numerical subjective probability.
\end{hbtheorem}

\begin{hbdefinition}{def:p17-ellsberg-three-colors}{Thirty red balls and sixty of unknown color}
An urn contains 30 red balls and 60 black or yellow balls
in an unknown split. One ball is drawn uniformly.
The four source acts, consistently in state order \((R,B,Y)\),
are
\[
\begin{array}{c|ccc}
 &R&B&Y\\ \hline
 F_1&H&L&L\\
 F_2&L&H&L\\
 F_3&H&L&H\\
 F_4&L&H&H
\end{array}
\]
The known information is \(p(R)=1/3\) and
\(p(B)+p(Y)=2/3\), rather than a known division between
black and yellow. For a fixed physical composition the
black probability is \(k/90\), \(k=0,\ldots,60\).
No uniform prior over these possible compositions is specified.
The principal pattern is
\(F_1\succ F_2,\ F_4\succ F_3\); the source also discusses
both strict comparisons reversed
\cite[pp.~653--656]{Ellsberg1961Ambiguity}.
\end{hbdefinition}

\begin{hbtheorem}{thm:p17-ellsberg-three-colors}{proved=annex:p17.ellsberg.three.colors}{The common-column contradiction}
Both strict patterns violate the sure-thing condition.
Neither can be represented by one additive prior and one
state-independent utility with \(u(H)>u(L)\).
For the principal pattern the two required inequalities
would be
\[
                          p(R)>p(B),\qquad p(R)<p(B).
\]
The contradiction holds without assuming a particular
curvature of utility in money. It is an exact representation
failure, with no implied numerical bound on an approximate
or stochastic model of observed choices.
\end{hbtheorem}

Ellsberg also considers small samples: one red and one black
from the unknown urn, or two of each; and one black and one
yellow from the group of sixty. The contradiction remains
conditional on the same preference pattern being retained.
Such a sample does not itself specify a posterior or a unique
set of plausible priors. A sampling and replacement protocol
is needed before calculating post-sampling composition
probabilities \cite[pp.~653--654]{Ellsberg1961Ambiguity}.

\section{What the original observations establish}

The original expressly describes the responses as gathered
under nonexperimental conditions and asks for controlled
experimentation to assess prediction. Its reports of more
frequent and less frequent patterns are not a participant count,
population percentage, or measured accuracy estimate.
The article describes people who retain the comparisons after
reflection, people who revise them to fit the axioms, and people
who remain uncertain. It also refers to Chipman's separate
experimental work; that work is a distinct source
\cite[p.~651, footnotes~7--9; pp.~655--656]{Ellsberg1961Ambiguity}.

The information distinction is also more specific than a contrast
between knowledge and total ignorance. The thirty-red-ball
information rules out many priors while leaving the black/yellow
division unsettled. Ellsberg discusses confidence in the quantity,
reliability, relevance and agreement of available evidence, and
allows a person to have a best probability estimate without treating
it as the sole decision criterion. His closing argument leaves the
normative dispute open; it does not prove that one alternative rule
is uniquely rational or delivers better long-run outcomes
\cite[pp.~657--663 and 669]{Ellsberg1961Ambiguity}.

\section{A specified finite set of priors}

\begin{hbdefinition}{def:p17-finite-maxmin}{The maxmin method and its inputs}
On the finite lottery-act domain, supply a nonempty finite
prior set \(\mathcal C=\{p^1,\ldots,p^m\}\subseteq\Delta(S)\),
\(m\geq1\), and the common utility \(u\). Define
\[
                 V_{\mathcal C}(f)=
                       \min_{p\in\mathcal C}E_p(f),\qquad
 f\succeq_{\mathcal C}g
                 \ \Longleftrightarrow\
                       V_{\mathcal C}(f)\geq V_{\mathcal C}(g).
\]
For a nonempty finite menu \(\mathcal M\) of available acts,
choose a maximizer of \(V_{\mathcal C}\) on \(\mathcal M\).
This is a worst-case \emph{expected-utility} rule over the
specified priors. It does not generally equal the minimum
utility over raw states.
The finite method is the specified-prior specialization of
the Gilboa--Schmeidler functional
\cite[pp.~141--145]{GilboaSchmeidler1989Maxmin}.
In its notation take \(X=Z\), \(\Sigma=2^S\),
\(Y=\Delta(Z)\), \(L_0=Y^S\), and the affine lottery utility
\(\widehat u(y)=\sum_z y(z)u(z)\).
The next theorem proves that the convex hull of the finite
prior list gives exactly the same functional.
Its prior set is an input to justify or compare, not an estimate
uniquely determined by a few urn choices.
\index{maxmin expected utility}\index{multiple priors}
\end{hbdefinition}

\begin{hbtheorem}{thm:p17-finite-maxmin-properties}{proved=annex:p17.finite.maxmin.properties}{Existence, continuity, convex hull and utility scale}
The minimum and every nonempty finite-menu maximum are attained.
The induced preference is complete and transitive, respects
coordinatewise improvement in \(x_f\), and, writing
\(\|x\|_\infty=\max_{s\in S}|x(s)|\), satisfies
\[
 |V_{\mathcal C}(f)-V_{\mathcal C}(g)|
                            \leq\|x_f-x_g\|_\infty,\qquad
 \min_s x_f(s)\leq V_{\mathcal C}(f)\leq\max_s x_f(s).
\]
Replacing \(\mathcal C\) by its convex hull changes no value.
Thus this finite model has the same functional as a closed
convex polytope of priors, with every minimum already attained
at a listed prior. Enlarging the finite prior list can only decrease each act's
value, although it can change the ranking of acts.
Replacing \(u\) by \(a u+b\), \(a>0\), changes every value
to \(aV_{\mathcal C}+b\) and preserves all rankings.
If \(u\) is constant, all acts are indifferent; if it is
nonconstant, some two constant acts are strictly ranked.
\end{hbtheorem}

\begin{hbdefinition}{def:p17-robust-dominance}{Comparison under every retained prior}
Write \(f\succeq_{\rm rob}g\) when
\[
                    E_p(f)\geq E_p(g)
                            \quad\hbox{for every }p\in\mathcal C.
\]
This robust dominance relation compares the same two acts
under every retained prior. It differs from comparing two
possibly different worst-case priors.
\end{hbdefinition}

\begin{hbtheorem}{thm:p17-robust-dominance}{proved=annex:p17.robust.dominance}{Dominance, incompleteness and the maxmin completion}
Robust dominance is reflexive and transitive and implies
\(f\succeq_{\mathcal C}g\). It need not compare every pair,
and maxmin preference does not imply robust dominance.
If \(E_p(f)>E_p(g)\) for every retained prior, then
\(V_{\mathcal C}(f)>V_{\mathcal C}(g)\).
Strict improvement under just one prior need not give this
strict maxmin conclusion.
For examples with consequence utilities \(0,1,2\),
two states and priors \(\delta_1,\delta_2\), utility
vectors \((1,0)\) and \((0,1)\) are robustly incomparable.
The maxmin rule strictly prefers \((1,1)\) to \((0,2)\),
although that comparison reverses under \(\delta_2\).
Vectors \((0,2)\) and \((0,1)\) illustrate weak robust
dominance with equal maxmin values.
\end{hbtheorem}

\begin{hbdefinition}{def:p17-ambiguity-mixtures}{Mixture and certainty independence}
For \(0\leq\lambda\leq1\), mix acts pointwise:
\[
           [\lambda f+(1-\lambda)g](s)
                        =\lambda f(s)+(1-\lambda)g(s).
\]
This is a roulette mixture conditional on the state, with
the randomizing mechanism independent of that state.
Full mixture independence requires, for every act \(h\)
and \(0<\lambda<1\),
\[
 f\succeq g\quad\Longleftrightarrow\quad
 \lambda f+(1-\lambda)h
                  \succeq\lambda g+(1-\lambda)h.
\]
Certainty independence restricts \(h\) to constant acts.
Uncertainty aversion in the source's axiom means
\[
 f\sim g\quad\Longrightarrow\quad
                       \lambda f+(1-\lambda)g\succeq f .
\]
These are the corresponding finite forms of
Gilboa--Schmeidler A.2 and A.5
\cite[pp.~144--145]{GilboaSchmeidler1989Maxmin}.
\end{hbdefinition}

\begin{hbtheorem}{thm:p17-maxmin-mixtures}{proved=annex:p17.maxmin.mixtures}{Concavity, certainty independence and the source's forward axioms}
The specified-prior value satisfies
\[
 V_{\mathcal C}(\lambda f+(1-\lambda)g)
           \geq\lambda V_{\mathcal C}(f)+(1-\lambda)V_{\mathcal C}(g).
\]
For a constant act \(h\), equality holds with its constant
utility \(x_h\):
\[
 V_{\mathcal C}(\lambda f+(1-\lambda)h)
                         =\lambda V_{\mathcal C}(f)+(1-\lambda)x_h.
\]
Hence certainty independence and uncertainty aversion hold.
The preference also has the source's mixture-continuity property:
if \(f\succ g\succ h\), there are \(\alpha,\beta\in(0,1)\)
such that
\[
       \alpha f+(1-\alpha)h\succ g
                              \succ\beta f+(1-\beta)h.
\]
Together with the preceding order and monotonicity results,
these verify the forward finite-model properties of A.1--A.5.
They do not infer a prior set from those axioms.
\end{hbtheorem}

\begin{hbtheorem}{thm:p17-maxmin-full-independence}{proved=annex:p17.maxmin.full.independence}{Exactly when full mixture independence survives}
On the entire finite lottery-act domain with nonconstant \(u\),
the maxmin preference satisfies full mixture independence
if and only if all priors in \(\mathcal C\) are identical.
With two states and priors \(\delta_1,\delta_2\),
the acts with normalized utility vectors
\(f=(1,0)\) and \(g=(0,1)\) are indifferent, while
\[
                V_{\mathcal C}(f)=0
                   <V_{\mathcal C}\bigl(\tfrac12f+\tfrac12g\bigr)=\tfrac12.
\]
This is both a hedging gain and a failure of full independence.
The nonconstant-utility and full-domain qualifications matter:
constant utility makes every comparison indifferent, and a
small observed menu need not expose every violated axiom.
\end{hbtheorem}

\section{Worked prior sets for the original comparisons}

\begin{hbtheorem}{thm:p17-maxmin-urns}{proved=annex:p17.maxmin.urns}{Finite worst-case calculations and nonidentification}
Normalize \(u(H)=1,u(L)=0\).
For the three-color information before sampling, use
\[
       \mathcal C_3=
           \{(1/3,0,2/3),\ (1/3,2/3,0)\}
                      \quad\hbox{in order }(R,B,Y).
\]
Their maxmin values are
\[
          (V(F_1),V(F_2),V(F_3),V(F_4))=(1/3,0,1/3,2/3).
\]
These minima are unchanged if all 61 physical compositions
or their convex hull are allowed. Thus the rule respects
the known red fraction and produces the principal strict pattern.

For the original two urns, use the auxiliary state pair
\((\hbox{color from I},\hbox{color from II})\), in order
\((RR,RB,BR,BB)\), and the finite prior list
\[
\begin{split}
 \mathcal C_2=\{&
 (1/2,1/2,0,0),\ (0,0,1/2,1/2),\\
 &(1/2,0,0,1/2),\ (0,1/2,1/2,0)\}.
\end{split}
\]
Its convex hull consists exactly of the joint distributions
with known-urn red marginal \(1/2\) and unrestricted
unknown-urn marginal. No independence between the auxiliary
colors is imposed. The four original single-urn bet values are
\[
          V(R_{\rm I})=V(B_{\rm I})=0,\qquad
          V(R_{\rm II})=V(B_{\rm II})=1/2 .
\]
For the marked 200-ball experiment, the prior endpoints
\((0,1/2,1/4,1/4)\) and \((1/2,0,1/4,1/4)\)
give values \(0,0,1/4,1/4,1/2,1/2\) for \(M_1,\ldots,M_6\).
They also give \(1/2,1/4,1/4,1/2\) for Arrow's four acts.

These observations do not identify one prior set.
For every \(0<d\leq1/3\), the two priors
\[
       (1/3,1/3-d,1/3+d),\quad
       (1/3,1/3+d,1/3-d)
\]
give the same strict three-color pattern, with values
\((1/3,1/3-d,2/3-d,2/3)\).
\end{hbtheorem}

\begin{hbtheorem}{thm:p17-ambiguity-randomization-averaging}{proved=annex:p17.ambiguity.randomization.averaging}{Adding a randomized act and averaging priors are different operations}
A fair independent coin used to select red or black before
betting on Urn I gives winning probability \(1/2\) under
every possible composition. Its maxmin value is \(1/2\)
on the normalized utility scale. It is an additional act
and can improve on either ambiguous-color bet alone.

By contrast, if fixed weights \(w_j\geq0\), \(\sum_jw_j=1\),
are used to average the prior expectations, then
\[
                \sum_jw_jE_{p^j}(f)=E_{\bar p}(f),
                           \qquad\bar p=\sum_jw_jp^j.
\]
That fixed averaging rule is one-prior expected utility
and satisfies the sure-thing condition. Averaging priors
therefore cannot produce the original strict reversal.
\end{hbtheorem}

\section{Ellsberg's own confidence-weighted rule}

\begin{hbdefinition}{def:p17-ellsberg-confidence-rule}{Estimated expectation and security level}
Supply a best-estimate prior \(p^0\), a nonempty finite
reasonable-prior set \(\mathcal C\), and confidence
\(\rho\in[0,1]\). Ellsberg's rule is
\[
           W_\rho(f)=\rho E_{p^0}(f)
                             +(1-\rho)\min_{p\in\mathcal C}E_p(f).
\]
The same inputs are held fixed across the comparison.
His \(\rho\) is a confidence parameter, not the probability
that a particular urn color occurs. This is the specific
proposed rule and numerical illustration on
pp.~664--668 \cite{Ellsberg1961Ambiguity}.
\end{hbdefinition}

\begin{hbtheorem}{thm:p17-ellsberg-confidence-rule}{proved=annex:p17.ellsberg.confidence.rule}{The exact original calculations and the finite maxmin bridge}
The confidence-weighted rule equals maxmin over the finite set
\[
                 \mathcal C_\rho=
                     \{\rho p^0+(1-\rho)p:p\in\mathcal C\}.
\]
For \(\mathcal C_3\), \(p^0=(1/3,1/3,1/3)\), and the
source's utility values \(u(H)=6,u(L)=0\), its values on
\(F_1,\ldots,F_4\) are
\[
                         (2,\ 2\rho,\ 2+2\rho,\ 4).
\]
At \(\rho=1/4\) they are \((2,1/2,5/2,4)\);
at \(\rho=3/4\) they are \((2,3/2,7/2,4)\).
Every \(\rho<1\) gives the principal strict reversal;
\(\rho=1\) gives the one-prior ties.
The effective worst-case priors for \(F_2,F_3\), respectively,
are
\[
       (1/3,\rho/3,(2-\rho)/3),\qquad
       (1/3,(2-\rho)/3,\rho/3).
\]
The prior set is common to the acts even though a minimizing
prior can depend on the act.
For the full finite consequence domain, put
\(D_u=\max_z u(z)-\min_z u(z)\). Then
\[
             |W_\rho(f)-E_{p^0}(f)|\leq(1-\rho)D_u
                         \quad\hbox{for every act }f.
\]
Thus the values converge uniformly to the best-estimate
expected utilities as \(\rho\to1\), even though the stated
strict urn reversals persist for every \(\rho<1\).

Replacing the minimum by a maximum, the optimistic variant
mentioned in the source gives
\((2,4-2\rho,6-2\rho,4)\), and hence the reversed strict
pattern for \(\rho<1\).
For a two-urn comparison with known red probability \(q\),
unknown-urn best estimate \(1/2\), and unknown-red security
level zero, the corresponding normalized values are
\(q\) and \(\rho/2\). Thus the known-red bet is preferred
exactly when \(q>\rho/2\), illustrating the source's
tradeoff as the known proportion becomes less favorable.
\end{hbtheorem}

All displayed weighted calculations retain the earlier acts
in state order \((R,B,Y)\). Exchanging the names of the two
unknown colors gives the source's later numerical presentation.
Confidence changes in this
illustration are specified comparisons, not a proved statistical
learning rule or a measured response curve.

\section{Other criteria and the limits of the comparison}

\begin{hbtheorem}{thm:p17-ambiguity-other-criteria}{proved=annex:p17.ambiguity.other.criteria}{Worst outcomes, Hurwicz values and the regret-menu qualification}
On the normalized three-color acts, each minimum state payoff
is zero and each maximum is one. Thus minimum-outcome,
maximum-outcome, and any fixed weighted average of these
two extremes tie the four acts. Pairwise minimax regret
also ties \(F_1,F_2\) and ties \(F_3,F_4\).
These raw-state criteria do not produce either strict source pattern.

For a fixed nonempty finite menu \(\mathcal M\) of utility vectors, define
\[
        {\rm Reg}_{\mathcal M}(x)
                 =\max_s\{\max_{y\in\mathcal M}y(s)-x(s)\}.
\]
Ranking by smaller regret is complete and transitive within
that menu, but need not satisfy the sure-thing condition.
For the fixed menu
\[
        \{(1,1),(0,1),(1,0),(0,0)\},
\]
the first vector is strictly preferred to the second,
whereas changing their common second coordinate to zero
produces two vectors with equal maximum regret.
Thus a claim about all regret-based decisions requires its
menu and benchmark convention; pairwise and fixed-menu
regret are not interchangeable.
\end{hbtheorem}

\begin{hbmethod}{def:p17-ambiguity-computation}{Using a finite ambiguity model}
Specify the states, common consequence utility, available acts,
known probability restrictions, and the finite prior list.
Compute each state utility \(x_f(s)\), then each \(E_{p^j}(f)\).
Retain all priorwise comparisons when robust dominance is the
question. For maxmin choice, record the minimum for each act
and select the largest minimum, including every tie.
If the confidence-weighted rule is chosen, also record
\(p^0,\rho\) and compute its stated index.
For finite rational inputs and a finite menu these are exact
finite sums and comparisons. Randomized mixtures are available
only when included in the decision domain.

Report the selected act, its priorwise expected utilities,
its worst-case value, and which input priors attain that value.
Changing the prior set, confidence weight, utility, or menu
can change the answer. None of the urn comparisons proves that
maxmin is uniquely rational, that the listed priors are uniquely
identified, or that a general multi-prior updating or representation
theorem has been established.
\end{hbmethod}
